% section_intro.tex
% §1: Introduction
% Author: Marvin L. Gentry

\section{Introduction}
\label{sec:intro}

\subsection{Background}

Two classical subjects of modern number theory — the arithmetic of
hyperbolic 3-manifolds and the theory of Bianchi modular forms — meet
in a single arithmetic object studied in this paper.

\medskip
\textbf{Hyperbolic 3-manifolds and their arithmetic.}
A \emph{closed hyperbolic 3-manifold} $M$ is a compact Riemannian
3-manifold of constant sectional curvature $-1$ and finite volume.
By the uniformisation theorem of Thurston, $M = \mathbb{H}^3 / \Gamma$
where $\Gamma \subset \mathrm{PSL}(2,\mathbb{C})$ is a discrete cocompact
torsion-free subgroup (the \emph{Kleinian group}).
The arithmetic of $M$ is encoded in three invariants of $\Gamma$:
\begin{itemize}
  \item The \emph{invariant trace field} $k(\Gamma)$, the subfield of
        $\mathbb{C}$ generated over $\mathbb{Q}$ by $\{\mathrm{tr}(\gamma)^2
        : \gamma \in \Gamma\}$ \cite{MFW};
  \item The \emph{invariant quaternion algebra} $A(\Gamma)$, a central
        simple algebra over $k(\Gamma)$ whose ramification controls the
        commensurability class of $\Gamma$;
  \item The \emph{character variety} $X(\Gamma, \mathrm{SL}(2,\mathbb{C}))$,
        the affine algebraic variety parameterising representations of
        $\pi_1(M)$ into $\mathrm{SL}(2,\mathbb{C})$ up to conjugacy.
\end{itemize}
For two-generator groups these invariants are largely captured by the
\emph{Fricke coordinates} $(x, y, z) = (\mathrm{tr}(a), \mathrm{tr}(b),
\mathrm{tr}(ab))$, which satisfy a cubic surface equation (the Fricke
identity) in the character variety \cite{Goldman}.

\medskip
\textbf{Bianchi modular forms.}
For an imaginary quadratic field $K = \mathbb{Q}(\sqrt{-d})$ with ring
of integers $\mathcal{O}_K$, the \emph{Bianchi group}
$\mathrm{PSL}(2, \mathcal{O}_K)$ acts on hyperbolic 3-space $\mathbb{H}^3$
by isometries. Its congruence subgroups and their associated spaces of
modular forms — \emph{Bianchi modular forms} — are the natural 3-dimensional
generalisation of classical modular forms over $\mathbb{Q}$.
The Langlands program predicts, and in many cases has established,
a correspondence between Bianchi newforms and Galois representations
\cite{TaylorWiles, Berger}.
In particular, Bianchi newforms over $\mathbb{Q}(\sqrt{-3})$ (the
Eisenstein field) are associated to elliptic curves over
$\mathbb{Q}(\sqrt{-3})$ and to arithmetic hyperbolic 3-manifolds fibered
over $\mathrm{Spec}\,\mathcal{O}_K$.

\subsection{The central object}

The hyperbolic 3-manifold at the centre of this paper is
\[
  M_0 = m006(-5,2),
\]
the closed manifold obtained by $(-5,2)$ Dehn surgery on the figure-eight
knot sister $m006$. It has volume $\approx 2.0289$, first homology
$H_1(M_0; \mathbb{Z}) = \mathbb{Z}/5\mathbb{Z}$, and is the 44th
manifold in the Hodgson--Weeks closed census ordered by volume \cite{snappy}.

The invariant trace field of $M_0$ is
\[
  K_{10} = \mathbb{Q}[t]/(p_{10}(t)),
  \quad
  p_{10}(t) = t^{10} - 7t^8 - 4t^7 + 17t^6 + 14t^5 - 18t^4 - 14t^3 + 8t^2 + 3t - 1,
\]
a degree-10 totally generic field (Galois group $S_{10}$) with discriminant
$\Delta(K_{10}) = -271488204251$ and signature $(r_1, r_2) = (8, 1)$.

On the automorphic side, there is a Bianchi newform $f_{283}$ of level
$\mathfrak{n}_{283}$ over $K = \mathbb{Q}(\sqrt{-3})$ whose Hecke
eigenvalues at rational primes agree with those of the elliptic curve
$E_1\colon y^2 + y = x^3 - x^2$ (Cremona label \texttt{11a1},
the modular curve $X_0(11)$). The form $f_{283}$ is associated to the
companion manifold
$M_{\mathrm{PMNS}} = m003(-2,3)$ (the Meyerhoff manifold, smallest-volume
closed hyperbolic 3-manifold with $H_1 = \mathbb{Z}/5$).

The two manifolds $M_0$ and $M_{\mathrm{PMNS}}$ share the homology group
$\mathbb{Z}/5\mathbb{Z}$ and the Bianchi base form $f_{283}$, but differ
in every other arithmetic invariant.

\subsection{Main results}

We establish three theorems that each describe a different facet of the
arithmetic of $M_0$ and its Bianchi form.

\medskip
\textbf{Theorem A (Frobenius discriminant, \S\ref{sec:frobenius}).}
Define the \emph{Farey tower primes} $p_k = 5k(k+1)+1$ ($p_1 = 11$,
$p_2 = 31$, $p_3 = 61$, \ldots). The $\chi_5^{(p_k)}$-twist of $f_{283}$
— where $\chi_5^{(p_k)}$ is the unique primitive order-5 character of
conductor $p_k$ — produces a $\mathbb{Q}(\sqrt{5})$-valued Bianchi newform
at level $283\cdot p_k$ if and only if the Frobenius discriminant satisfies
\[
  a_{p_k}(f_{11})^2 - 4p_k = -3\,m^2 \quad\text{for some }m \in \mathbb{Z}.
\]
Among all tower primes $p_k \leq 281$, this condition holds for exactly
one prime: $p_2 = 31$, where $7^2 - 4\cdot 31 = -75 = -3\cdot 5^2$.
The $\mathbb{Q}(\sqrt{5})$ Bianchi component at level $8773 = 283\cdot 31$
is confirmed in the LMFDB.

\medskip
\textbf{Theorem B (Three-Ray Theorem, \S\ref{sec:three-ray}).}
The Hecke eigenvalues of the $\mathbb{Q}(\sqrt{5})$ Bianchi newform $g$
at level 8773 satisfy
\[
  a_p(g) = a_p(f_{11}) \cdot S_{k(p)},
  \qquad
  k(p) = \mathrm{dlog}_3(p \bmod 31) \bmod 5,
\]
where $S_k = \zeta_5^k + \zeta_5^{-k}$ takes exactly three values:
$2$, $1/\varphi$, and $-\varphi$ (with $\varphi$ the golden ratio).
Every Hecke eigenvalue of $g$ therefore lies on one of three rays through
the origin in $\mathbb{Q}(\sqrt{5}) \cong \mathbb{R}^2$.
The coefficient field $\mathbb{Q}(\sqrt{5})$ is generated, as a
$\mathbb{Q}$-algebra, by the non-trivial Hecke multipliers
$\{1/\varphi, -\varphi\}$ of $g$ itself — a self-referential structure
specific to the case of fifth roots of unity.

\medskip
\textbf{Theorem C (Trace Quotient Theorem, \S\ref{sec:trace-quotient}).}
The Dehn surgery relation of $M_0 = m006(-5,2)$ forces the Fricke
coordinate collapse
\[
  \mathrm{tr}(\rho(ab)) = \mathrm{tr}(\rho(a)),
\]
reducing the character variety from a cubic surface to a two-variable
surface. This collapse propagates to a cascade of exact trace class
equalities, and implies that the search space of words of length $\leq 6$
maps to a finite trace quotient graph with at most 122 nodes. Furthermore,
the generator of $K_{10}$ is identified as
\[
  \mathrm{tr}(\rho(a)) = -\alpha,
\]
where $\alpha$ is the canonical root of $p_{10}$.

Among all 948 closed hyperbolic 3-manifolds with $H_1 = \mathbb{Z}/5$
in the SnapPy closed census (11,031 manifolds by volume), $M_0$ is the
\emph{unique} one whose invariant trace field generator has signature
$(8,1)$. All others have $r_2 \geq 2$.

\subsection{The common thread}

Theorems A, B, and C arise from a single arithmetic coincidence: the
Hecke eigenvalue $a_{31}(f_{11}) = 7$ satisfies $7^2 - 4\cdot 31 = -3\cdot5^2$.
This identity:
\begin{enumerate}[(i)]
  \item forces $p = 31$ as the unique tower prime admitting Bianchi descent
        (Theorem A);
  \item forces the eigenvalue field $\mathbb{Q}(\sqrt{5})$, generated by
        $\zeta_5 + \zeta_5^{-1} = 1/\varphi$, whose generator is the
        Hecke multiplier (Theorem B);
  \item is connected to $M_0$ through its surgery coefficient $(-5,2)$,
        which forces $H_1 = \mathbb{Z}/5$ and the Fricke collapse
        $z = x$ (Theorem C).
\end{enumerate}

The chain $31 \to \chi_5 \to \mathbb{Q}(\sqrt{5}) \to \varphi \to M_0$
is thus forced by the single numerical fact $a_{31} = 7$.

\subsection{Conventions and notation}

Throughout, we write:
\begin{itemize}
  \item $K = \mathbb{Q}(\sqrt{-3})$, with ring of integers
        $\mathcal{O}_K = \mathbb{Z}[\omega]$, $\omega = e^{2\pi i/3}$;
  \item $f_{11}$: the weight-2 newform of level 11
        (Cremona label \texttt{11a1}), with $a_p = a_p(f_{11})$;
  \item $f_{283}$: the associated Bianchi newform over $K$ at level $\mathfrak{n}_{283}$;
  \item $g$: the $\mathbb{Q}(\sqrt{5})$-Bianchi newform at level 8773;
  \item $M_0 = m006(-5,2)$: the central manifold, with
        $\rho\colon \pi_1(M_0) \to \mathrm{SL}(2,\mathbb{C})$
        its polished holonomy;
  \item $\alpha$: the root of $p_{10}$ with $\mathrm{Im}(\alpha) < 0$,
        so that $\mathrm{tr}(\rho(a)) = -\alpha$;
  \item $\varphi = (1+\sqrt{5})/2$: the golden ratio.
\end{itemize}
We use $A = a^{-1}$, $B = b^{-1}$ for the inverses of the standard
generators.

\subsection{Computational methods and reproducibility}

All results are verified computationally in SageMath (version 10.x)
and SnapPy (version 3.x). The verification scripts are available in
the \texttt{reproduce/} directory of the accompanying repository:
\begin{itemize}
  \item \texttt{verify\_frobenius.py}: reproduces Table~\ref{tab:frobenius}
        (Frobenius discriminant at each tower prime);
  \item \texttt{verify\_three\_ray.py}: verifies the three-ray structure
        for all primes $p < 300$;
  \item \texttt{verify\_trace.py}: verifies the Fricke collapse, trace
        class equalities, and ITF generator identity;
  \item \texttt{signature\_enum\_test.py}: exhaustive enumeration of
        ITF signatures among all 948 $H_1 = \mathbb{Z}/5$ manifolds
        in the closed census.
\end{itemize}
Each script prints \texttt{PASS}/\texttt{FAIL} for each claimed equality
and can be run in under 10 minutes on a standard laptop.

\subsection{Organisation}

Section~\ref{sec:frobenius} proves Theorem A and establishes the
Frobenius discriminant criterion that forces $p = 31$.
Section~\ref{sec:three-ray} proves Theorem B and develops the
three-ray and self-referential structures of the Hecke eigenvalues.
Section~\ref{sec:trace-quotient} proves Theorem C, including the
ITF generator identity and the uniqueness of the $(8,1)$ signature.
Section~\ref{sec:connections} discusses connections among the three
theorems and states the remaining open questions.

\medskip
\noindent\textbf{Acknowledgements.}
Computations were performed using SnapPy \cite{snappy} and SageMath.
The LMFDB \cite{lmfdb} database was essential for independent verification
of the Bianchi form at level 8773.
